\documentclass[handout]{beamer}
\mode<presentation>
\usetheme[secheader]{Boadilla}
\usecolortheme{whale}

\usepackage[utf8]{inputenc}
\usepackage[T1]{fontenc}
\usepackage[indonesian]{babel}
\usepackage{lmodern}
\usepackage{graphicx}
\usepackage{bookmark}

\setbeamertemplate{footline}
{
	\leavevmode%
	\hbox{%
		\begin{beamercolorbox}[wd=.333333\paperwidth,ht=2.25ex,dp=1ex,center]{author in head/foot}%
			\usebeamerfont{author in head/foot}\insertshortauthor%
		\end{beamercolorbox}%
		\begin{beamercolorbox}[wd=.433333\paperwidth,ht=2.25ex,dp=1ex,center]{title in head/foot}%
			\usebeamerfont{title in head/foot}\insertshorttitle
		\end{beamercolorbox}%
		\begin{beamercolorbox}[wd=.233333\paperwidth,ht=2.25ex,dp=1ex,right]{date in head/foot}%
			\usebeamerfont{date in head/foot}\insertshortdate{}\hspace*{2em} \insertframenumber{} / \inserttotalframenumber\hspace*{2ex}
	\end{beamercolorbox}}%
	\vskip0pt%
}

\title[Bab 9 -- JavaScript \& DOM]{Pemrograman Web}
\subtitle{Bab 9: JavaScript Dasar, DOM, dan Manipulasi Elemen}
\author{Cecep Suwanda, S.Si., M.Kom.}
\institute{Universitas Bale Bandung}
\date{\today}

\begin{document}

% Slide Judul
\begin{frame}
	\titlepage
\end{frame}

% Outline dan CPMK
\begin{frame}{Outline Bab 9}
	\begin{block}{Sub-CPMK 2.3}
		Mengimplementasikan interaktivitas dengan JavaScript (DOM, event handling, validasi form).
	\end{block}
	\begin{itemize}
		\item Pengenalan JavaScript dan Perannya di Web
		\item Menjalankan JavaScript di Browser
		\item Sintaks Dasar dan Variabel (var, let, const)
		\item Tipe Data dan Operator
		\item Konversi Tipe, String
		\item Percabangan dan Perulangan
		\item Fungsi, Scope, Closure
		\item Array dan Objek
		\item Pengenalan DOM
		\item Mengakses Elemen dan Mengubah Konten
		\item Event Handling dan Validasi Form
	\end{itemize}
\end{frame}

% ========== SECTION 9-1: Pengenalan JavaScript ==========
\begin{frame}{Pengenalan JavaScript}
	\begin{itemize}
		\item JavaScript: bahasa pemrograman di sisi \textbf{client} (browser); menambah \textbf{perilaku}---interaktivitas.
		\item HTML = struktur, CSS = tampilan, JavaScript = perilaku (sistem listrik).
		\item Diciptakan Brendan Eich (Netscape, 1995); standar \textbf{ECMAScript}.
		\item ES6/ES2015: \texttt{let}/\texttt{const}, arrow function, class, template literal, Promise.
		\item Ekosistem: browser, Node.js (server), React Native (mobile), Electron (desktop).
	\end{itemize}
	\begin{block}{Catatan}
		JavaScript $\neq$ Java. JS dinamis, interpreted, prototipe; Java statically typed, compiled, berbasis class.
	\end{block}
\end{frame}

% ========== SECTION 9-2: Menjalankan JavaScript ==========
\begin{frame}{Menjalankan JavaScript di Browser}
	\begin{itemize}
		\item \textbf{Inline}: \texttt{onclick="alert('Hi')"} --- hindari.
		\item \textbf{Internal}: \texttt{<script>} di HTML --- untuk demo kecil.
		\item \textbf{Eksternal}: \texttt{<script src="app.js">} --- \textbf{direkomendasikan}.
		\item \textbf{DevTools Console}: F12 atau Ctrl+Shift+I; \texttt{console.log()}, \texttt{console.error()}, \texttt{console.table()}.
		\item Atribut \texttt{defer}: tunda eksekusi sampai HTML selesai di-parse (urutan terjaga).
		\item Atribut \texttt{async}: jalankan segera setelah diunduh (urutan tidak dijamin). Untuk DOM: gunakan \texttt{defer}.
	\end{itemize}
\end{frame}

% ========== SECTION 9-3, 9-4: Sintaks dan Variabel ==========
\begin{frame}{Sintaks Dasar dan Variabel}
	\begin{itemize}
		\item Statement diakhiri \texttt{;} (ASI ada, tapi tulis eksplisit).
		\item Case-sensitive; camelCase, PascalCase, UPPER\_SNAKE\_CASE.
		\item \texttt{"use strict";} --- mode ketat, lebih aman.
	\end{itemize}
	\begin{block}{var vs let vs const}
		\texttt{var}: function scope, hoisted, re-declare OK. \\
		\texttt{let}: block scope, re-assign OK. \texttt{const}: block scope, tidak re-assign. \\
		\textbf{Rekomendasi}: default \texttt{const}; \texttt{let} jika perlu ubah; hindari \texttt{var}.
	\end{block}
\end{frame}

% ========== SECTION 9-5, 9-6: Tipe Data dan Operator ==========
\begin{frame}{Tipe Data dan Operator}
	\begin{itemize}
		\item \textbf{Primitif}: number, string, boolean, undefined, null. \texttt{typeof} untuk cek tipe.
		\item \textbf{Truthy/Falsy}: falsy = false, 0, "", null, undefined, NaN.
		\item \textbf{Operator}: aritmetika (+, -, *, /, \%, **), perbandingan (===, !==), logika (\&\&, ||, !).
		\item \textbf{Selalu gunakan \texttt{===}} (strict equality); hindari \texttt{==} (type coercion).
		\item \texttt{??} (nullish coalescing): ganti hanya jika null/undefined.
	\end{itemize}
\end{frame}

% ========== SECTION 9-7, 9-8: Konversi dan String ==========
\begin{frame}{Konversi Tipe dan String}
	\begin{itemize}
		\item \textbf{Type coercion}: \texttt{"5" + 3} $\to$ \texttt{"53"}; \texttt{"5" - 3} $\to$ 2.
		\item \texttt{Number(x)}, \texttt{parseInt()}, \texttt{parseFloat()}; \texttt{String(x)}; \texttt{Boolean(x)}.
		\item \texttt{NaN !== NaN}; gunakan \texttt{Number.isNaN(x)}.
		\item \textbf{String}: \texttt{.length}, \texttt{.toUpperCase()}, \texttt{.includes()}, \texttt{.slice()}, \texttt{.split()}, \texttt{.trim()}.
		\item \textbf{Template literal}: \texttt{`Hello \$\{name\}!`} --- interpolasi, multi-baris.
	\end{itemize}
\end{frame}

% ========== SECTION 9-9, 9-10: Percabangan dan Perulangan ==========
\begin{frame}{Percabangan dan Perulangan}
	\begin{itemize}
		\item \texttt{if}, \texttt{else if}, \texttt{else}; \texttt{switch} (jangan lupa \texttt{break}).
		\item Ternary: \texttt{kondisi ? nilaiTrue : nilaiFalse}.
		\item \texttt{for}, \texttt{while}, \texttt{do...while}; \texttt{break}, \texttt{continue}.
		\item \texttt{for...of}: iterasi elemen array; \texttt{for...in}: iterasi key objek (hindari untuk array).
	\end{itemize}
\end{frame}

% ========== SECTION 9-11, 9-12: Fungsi dan Scope ==========
\begin{frame}{Fungsi dan Scope}
	\begin{itemize}
		\item \textbf{Tiga cara}: function declaration (hoisted), function expression, arrow function.
		\item Default parameter: \texttt{function f(a, b = 2) \{\}}.
		\item Arrow function: lexical \texttt{this}; cocok untuk callback; tidak untuk method objek.
		\item \textbf{Scope}: global, function (\texttt{var}), block (\texttt{let}/\texttt{const}).
		\item \textbf{Closure}: fungsi ``mengingat'' variabel scope luar. Hindari variabel global.
	\end{itemize}
\end{frame}

% ========== SECTION 9-13, 9-14: Array dan Objek ==========
\begin{frame}{Array}
	\begin{itemize}
		\item Indeks dari 0; dinamis; literal \texttt{[]}.
		\item \texttt{push}, \texttt{pop}, \texttt{unshift}, \texttt{shift}, \texttt{slice}, \texttt{splice}.
		\item \texttt{map}, \texttt{filter}, \texttt{reduce} --- metode fungsional.
		\item Destructuring: \texttt{const [a, b, ...rest] = arr;}
	\end{itemize}
\end{frame}

\begin{frame}{Objek}
	\begin{itemize}
		\item Literal \texttt{\{\}}; key-value; dot notation \texttt{obj.key} atau bracket \texttt{obj["key"]}.
		\item Properti + metode (fungsi di dalam objek).
		\item \texttt{Object.keys(obj)}; \texttt{"key" in obj}; spread \texttt{\{...obj\}}.
		\item Destructuring: \texttt{const \{nama, nim\} = mahasiswa;}
		\item Array untuk koleksi terurut; Objek untuk entitas dengan properti bernama.
	\end{itemize}
\end{frame}

% ========== SECTION 9-15: DOM ==========
\begin{frame}{Pengenalan DOM}
	\begin{itemize}
		\item \textbf{DOM} (Document Object Model): representasi HTML sebagai pohon objek.
		\item Setiap elemen = \textbf{node}; \texttt{document} = titik masuk.
		\item DOM adalah \textit{live} --- mencerminkan keadaan saat ini, bukan file HTML asli.
		\item Metode akses: \texttt{getElementById()}, \texttt{querySelector()}, \texttt{querySelectorAll()}.
	\end{itemize}
\end{frame}

% ========== SECTION 9-16: Mengakses dan Mengubah DOM ==========
\begin{frame}{Mengakses Elemen dan Mengubah Konten}
	\begin{itemize}
		\item \texttt{querySelector(".css")} --- 1 elemen; \texttt{querySelectorAll("p")} --- NodeList.
		\item \texttt{textContent}: teks (aman XSS); \texttt{innerHTML}: HTML (risiko XSS dari input pengguna).
		\item \texttt{setAttribute()}, \texttt{el.style}, \texttt{el.classList.add/toggle/remove}.
		\item Untuk input pengguna: gunakan \texttt{textContent} atau \texttt{createElement()}+\texttt{appendChild()}.
	\end{itemize}
\end{frame}

% ========== SECTION 9-17: Event ==========
\begin{frame}{Event Handling}
	\begin{itemize}
		\item \texttt{addEventListener("event", handler)} --- click, submit, load, input, keydown, change.
		\item \texttt{event.preventDefault()} --- cegah perilaku default (mis. submit form).
		\item Validasi form: event \texttt{submit} + \texttt{preventDefault()} jika gagal.
		\item \textbf{Event bubbling}: event naik dari target ke parent; \texttt{event.target} vs \texttt{event.currentTarget}.
		\item \textbf{Event delegation}: satu listener di parent untuk banyak child --- efisien untuk daftar dinamis.
	\end{itemize}
\end{frame}

% ========== Rangkuman dan Penutup ==========
\begin{frame}{Rangkuman Bab 9}
	\begin{itemize}
		\item \texttt{const} default; \texttt{let} jika re-assign; hindari \texttt{var}.
		\item Selalu \texttt{===}; waspadai type coercion.
		\item Arrow function untuk callback; function declaration untuk method.
		\item DOM: \texttt{querySelector}, \texttt{textContent} (bukan \texttt{innerHTML}) untuk keamanan.
		\item Event: \texttt{addEventListener}; event delegation untuk efisiensi.
	\end{itemize}
\end{frame}

\begin{frame}{}
	\centering
	\vfill
	\textbf{\Large Pertanyaan?}
	\vfill
\end{frame}

\end{document}
